% ===== PRÉAMBULE CANONIQUE — à coller VERBATIM en tête de CHAQUE .tex =====
\documentclass[11pt,a4paper]{article}
\usepackage[utf8]{inputenc}\usepackage[T1]{fontenc}\usepackage{lmodern}
\usepackage[french]{babel}
\usepackage{amsmath,amssymb,mathtools}\usepackage{icomma}
\usepackage[version=4]{mhchem}          % \ce{...} : équations & formules
\usepackage{siunitx}\sisetup{locale=FR} % \qty{}{}, \si{}, \ang{}
\usepackage{chemfig}                    % \chemfig{...} : structures développées
\usepackage{xcolor}\usepackage{colortbl}
\usepackage{tikz}\usepackage{pgfplots}\pgfplotsset{compat=1.18}
\usepackage[most]{tcolorbox}\usepackage{enumitem}\usepackage{array,booktabs}
\usepackage[a4paper,margin=2.2cm]{geometry}
\usetikzlibrary{arrows.meta,calc,angles,quotes,positioning,patterns,backgrounds,shapes.geometric}
\definecolor{primary}{RGB}{222,49,99}\definecolor{primarydark}{RGB}{150,22,55}
\definecolor{defcol}{RGB}{0,114,178}\definecolor{methcol}{RGB}{52,92,148}
\definecolor{excol}{RGB}{0,135,90}\definecolor{attncol}{RGB}{200,40,55}
\tcbset{boxbase/.style={enhanced,breakable,boxrule=0.9pt,arc=2.5pt,
  left=3mm,right=3mm,top=2.3mm,bottom=2.3mm,fonttitle=\bfseries,coltitle=white}}
% --- boîtes TUTORIEL (noms EXACTS, ne pas renommer) ---
\newtcolorbox{cle}[1][]{boxbase,colback=defcol!6,colframe=defcol,
  title={\textsc{À retenir}\ifx\relax#1\relax\relax\else\textmd{ : \itshape #1}\fi}}
\newtcolorbox{methode}[1][]{boxbase,colback=methcol!7,colframe=methcol,sharp corners,
  title={\textsc{Méthode}\ifx\relax#1\relax\relax\else\textmd{ --- \itshape #1}\fi}}
\newtcolorbox{exemplebox}[1][]{boxbase,colback=excol!5,colframe=excol,
  title={\textsc{Exemple}\ifx\relax#1\relax\relax\else\textmd{ : \itshape #1}\fi}}
\newtcolorbox{piege}[1][]{boxbase,colback=attncol!6,colframe=attncol,
  title={\textsc{Piège}\ifx\relax#1\relax\relax\else\textmd{ : \itshape #1}\fi}}
\newtcolorbox{remarque}[1][]{boxbase,colback=black!4,colframe=black!45,
  title={\textsc{Remarque}\ifx\relax#1\relax\relax\else\textmd{ : \itshape #1}\fi}}
% --- boîtes EXERCICE / CORRIGÉ (noms EXACTS, auto-comptées) ---
\newtcolorbox[auto counter]{exo}[1][]{boxbase,colback=primary!4,colframe=primary,
  title={\textsc{Exercice}~\thetcbcounter\ifx\relax#1\relax\relax\else\textmd{ --- \itshape #1}\fi}}
\newtcolorbox[auto counter]{corrige}[1][]{boxbase,colback=excol!4,colframe=excol,
  title={\textsc{Corrigé}~\thetcbcounter\ifx\relax#1\relax\relax\else\textmd{ --- \itshape #1}\fi}}
\newcommand{\R}{\mathbb{R}}\newcommand{\N}{\mathbb{N}}\newcommand{\Z}{\mathbb{Z}}\newcommand{\ds}{\displaystyle}
% (un \usepackage{fancyhdr}+en-tête est possible ; il est ignoré côté web)
% =========================================================================
\begin{document}
% THEME_TITLE: Formule de Nernst
\begin{center}\color{primarydark}\Large\bfseries Thème 4 — La formule de Nernst\end{center}

En conditions standard, un couple a le potentiel $E_{\mathrm{r}}^{\circ}$. \textbf{Hors} de ces
conditions, le potentiel dépend des \textbf{concentrations} : c'est la \textbf{formule de Nernst}.

\begin{cle}[formule de Nernst à $25\,^{\circ}\mathrm{C}$]
Pour une réaction redox échangeant $n$ électrons :
\[
\Delta E=\Delta E^{\circ}-\frac{0{,}06}{n}\log Q
\qquad\Bigl(\text{forme exacte : }\Delta E^{\circ}-\tfrac{0{,}0592}{n}\log Q\Bigr)
\]
où $Q$ est le \textbf{quotient réactionnel}. Le facteur $\tfrac{RT}{F}\ln 10=0{,}0592\approx 0{,}06$~V à
$25\,^{\circ}$C.
\end{cle}

\begin{cle}[potentiel d'une demi-pile]
Pour une demi-équation $\mathrm{Ox}+n\,\ce{e^-}\rightleftharpoons\mathrm{Red}$ :
\[
E=E_{\mathrm{r}}^{\circ}+\frac{0{,}06}{n}\log\frac{[\mathrm{Ox}]}{[\mathrm{Red}]}
\]
Les \textbf{solides} et l'\textbf{eau} (solvant) ont une activité $=1$ : ils n'apparaissent
\emph{pas} dans le rapport. Si la demi-équation contient des \ce{H+}, ceux-ci figurent dans $[\mathrm{Ox}]$
avec leur \textbf{coefficient en exposant} (le potentiel dépend alors du \textbf{pH}).
\end{cle}

\begin{methode}[calculer le potentiel d'une demi-pile]
\begin{enumerate}[label=\textbf{\arabic*.},leftmargin=2.0em,topsep=1pt,itemsep=1pt]
  \item Écrire la demi-équation et repérer $n$.
  \item Former le rapport $[\mathrm{Ox}]/[\mathrm{Red}]$ (solide et eau $\to 1$).
  \item Reporter dans $E=E_{\mathrm{r}}^{\circ}+\tfrac{0{,}06}{n}\log\frac{[\mathrm{Ox}]}{[\mathrm{Red}]}$.
\end{enumerate}
\end{methode}

\begin{exemplebox}[électrode de cuivre diluée]
\ce{Cu^{2+} + 2 e^- <=> Cu}, $n=2$, avec $[\ce{Cu^{2+}}]=0{,}01$~mol/L et
$E_{\mathrm{r}}^{\circ}=+0{,}34$~V (le métal \ce{Cu} a une activité $1$) :
\[
E=0{,}34+\frac{0{,}06}{2}\log(0{,}01)=0{,}34+0{,}03\times(-2)=0{,}34-0{,}06=+0{,}28~\text{V}.
\]
Diluer la forme oxydée \textbf{abaisse} le potentiel.
\end{exemplebox}

\begin{cle}[force électromotrice d'une pile]
\[
\Delta E=E_{\text{cathode}}-E_{\text{anode}},\qquad
\Delta E^{\circ}=E_{\mathrm{r}}^{\circ}(\text{cathode})-E_{\mathrm{r}}^{\circ}(\text{anode}).
\]
On calcule le potentiel de chaque demi-pile par Nernst, puis on fait la différence
(cathode $=$ couple le plus haut). À l'\textbf{équilibre}, $\Delta E=0$.
\end{cle}

\begin{piege}[réflexes de calcul]
\begin{itemize}[leftmargin=1.5em,topsep=1pt,itemsep=0pt]
  \item $\log(1)=0$ (on retrouve $E=E_{\mathrm{r}}^{\circ}$), $\log(0{,}1)=-1$, $\log(0{,}01)=-2$,
        $\log(10)=+1$.
  \item Ne pas oublier le $\tfrac{1}{n}$ : le facteur est $\tfrac{0{,}06}{n}$, \emph{pas} $0{,}06$.
  \item Une demi-équation avec \ce{H+} dépend fortement du pH ($[\ce{H+}]$ à la puissance de son
        coefficient).
\end{itemize}
\end{piege}

\end{document}
